\section*{अध्याय \ref{ch:g5:water-shapes-the-land} --- पानी ज़मीन को आकार देता है}

\begin{solution}{exo:g5:water-shapes-the-land:1}
(क) \omterm{def:g5:water-shapes-the-land:weathering}{अपक्षय}: \omterm{def:g3:stones-pebbles-sand:rock}{चट्टान} अपनी जगह पर टूटती है। (ख) \omterm{def:g5:water-shapes-the-land:erosion}{अपरदन}: \omterm{def:g3:stones-pebbles-sand:sand}{रेत} घिसकर बहा ले जाई
जाती है। (ग) \omterm{def:g5:water-shapes-the-land:deposition}{निक्षेपण}: \omterm{def:g3:stones-pebbles-sand:sand}{कीचड़} नीचे बैठता है। (घ)
\omterm{def:g5:water-shapes-the-land:weathering}{अपक्षय}: जड़ पत्थर को उसकी जगह पर ही तोड़ती है।
\end{solution}

\begin{solution}{exo:g5:water-shapes-the-land:2}
\omterm{def:g5:water-shapes-the-land:source-mouth}{उद्गम} वह जगह है जहाँ से \omterm{def:g1:landscapes-around-us:water}{नदी} शुरू होती है (एक सोता, पिघलता हिमनद, कोई ऊँची
गीली ज़मीन)। \omterm{def:g5:water-shapes-the-land:source-mouth}{मुहाना} वह जगह है जहाँ वह ख़त्म होती है, \omterm{def:g2:water-on-earth:ocean}{समुद्र}, किसी \omterm{def:g1:landscapes-around-us:water}{झील} या किसी
बड़ी \omterm{def:g1:landscapes-around-us:water}{नदी} में। \omterm{def:g5:water-shapes-the-land:source-mouth}{सहायक नदी} एक छोटी \omterm{def:g1:landscapes-around-us:water}{नदी} है जो किसी बड़ी \omterm{def:g1:landscapes-around-us:water}{नदी} में आकर
मिलती है।
\end{solution}

\begin{solution}{exo:g5:water-shapes-the-land:3}
\omterm{def:g5:water-shapes-the-land:source-mouth}{उद्गम}; सँकरी V-आकार की \omterm{def:g1:landscapes-around-us:landform}{घाटी}; \omterm{def:g5:water-shapes-the-land:meander}{विसर्प}; \omterm{def:g5:water-shapes-the-land:delta}{डेल्टा}।
\end{solution}

\begin{solution}{exo:g5:water-shapes-the-land:4}
(क) ग़लत: पाला, पानी, जड़ें, गर्मी और ठंड कठोर \omterm{def:g3:stones-pebbles-sand:rock}{चट्टान} को भी धीरे-धीरे तोड़ते हैं
(\omterm{def:g5:water-shapes-the-land:weathering}{अपक्षय})। (ख) सही: पानी जितना तेज़, उतने बड़े टुकड़े वह हिला सकता है। (ग) ग़लत:
सबसे भारी टुकड़े पहले बैठते हैं; सबसे बारीक \omterm{def:g3:stones-pebbles-sand:sand}{कीचड़} सबसे आख़िर में बैठता है, जहाँ
पानी लगभग ठहरा हो, किसी \omterm{def:g1:landscapes-around-us:water}{झील} में या \omterm{def:g2:water-on-earth:ocean}{समुद्र} में।
\end{solution}

\begin{solution}{exo:g5:water-shapes-the-land:5}
A मोड़ के बाहरी ओर है: वहाँ पानी तेज़ और गहरा है, और वह किनारे को घिसता है
(\omterm{def:g5:water-shapes-the-land:erosion}{अपरदन})। B मोड़ के भीतरी ओर है: वहाँ पानी धीमा है, और वह \omterm{def:g3:stones-pebbles-sand:sand}{रेत} गिराता है
(\omterm{def:g5:water-shapes-the-land:deposition}{निक्षेपण})।
\end{solution}

\begin{solution}{exo:g5:water-shapes-the-land:6}
पहाड़ों में \omterm{def:g1:landscapes-around-us:water}{नदी} खड़ी ढलान वाली और तेज़ है: वह बड़े पत्थरों को लुढ़का सकती है, और
छोटे कणों को बहा ले जाती है। \omterm{def:g2:water-on-earth:ocean}{समुद्र} के पास \omterm{def:g1:landscapes-around-us:water}{नदी} धीमी है: वह अब बड़े पत्थर नहीं
हिला सकती, केवल \omterm{def:g3:stones-pebbles-sand:sand}{रेत} और \omterm{def:g3:stones-pebbles-sand:sand}{कीचड़}, और उन्हें वहीं गिरा देती है। रास्ते में पत्थर
टूटकर और घिसकर छोटे कण भी बन गए
हैं।
\end{solution}

\begin{solution}{exo:g5:water-shapes-the-land:7}
(क) ऊपरी मार्ग: $600 \div 6 = 100$ किलोमीटर। मध्य मार्ग:
$600 \div 2 = 300$ किलोमीटर। (ख) निचला मार्ग: $600 - 100 - 300 = 200$
किलोमीटर। क्योंकि $3 \times 200 = 600$, वह \omterm{def:g1:landscapes-around-us:water}{नदी} का $\frac{1}{3}$ है।
(जाँच: $\frac{1}{6} + \frac{1}{2} + \frac{1}{3} = \frac{1}{6} +
\frac{3}{6} + \frac{2}{6} = \frac{6}{6}$।)
\end{solution}

\begin{solution}{exo:g5:water-shapes-the-land:8}
हल्के रंग के धब्बे मोड़ों के भीतरी ओर बिछी \omterm{def:g3:stones-pebbles-sand:sand}{रेत} और \omterm{def:g3:stones-pebbles-sand:pebble}{बजरी} हैं, जहाँ पानी धीमा है
(\omterm{def:g5:water-shapes-the-land:deposition}{निक्षेपण})। गहरे रंग की घुमावदार \omterm{def:g1:landscapes-around-us:water}{झीलें} \omterm{def:g1:landscapes-around-us:water}{नदी} के पुराने फंदे हैं: दो फंदे इतने पास आ
गए कि \omterm{def:g1:landscapes-around-us:water}{नदी} उनके बीच की ज़मीन की गर्दन को तोड़कर छोटे रास्ते से निकल गई, और पुराना
फंदा पीछे एक \omterm{def:g1:landscapes-around-us:water}{झील} बनकर
रह गया।
\end{solution}

\begin{solution}{exo:g5:water-shapes-the-land:9}
\omterm{def:g1:landscapes-around-us:water}{नदी} के बगल के समतल चरागाह उसका \omterm{def:g5:water-shapes-the-land:flood}{बाढ़ का मैदान} हैं: वह ज़मीन जिसे \omterm{def:g1:landscapes-around-us:water}{नदी} \omterm{def:g5:water-shapes-the-land:flood}{बाढ़} आने पर
ढक लेती है। वे ठीक इसीलिए समतल हैं कि \omterm{def:g1:landscapes-around-us:water}{नदी} ने उन्हें अपने ही \omterm{def:g3:stones-pebbles-sand:sand}{कीचड़} से बनाया है।
किसी दिन \omterm{def:g1:landscapes-around-us:water}{नदी} में फिर \omterm{def:g5:water-shapes-the-land:flood}{बाढ़} आएगी और स्कूल में पानी भर देगी, और \omterm{def:g5:water-shapes-the-land:flood}{बाढ़} का तेज़ पानी
ख़तरनाक है। स्कूल ऊँची ज़मीन पर बनना चाहिए, और चरागाहों को खेतों, चरागाहों या
एक उद्यान के रूप में रखना चाहिए।
\end{solution}

\begin{solution}{exo:g5:water-shapes-the-land:10}
(क) एक साल में 12 महीने होते हैं: $12$ अरब $\div\, 12 = 1$ अरब टन हर महीने,
औसतन। (ख) एक अरब $1\,000\,000\,000$ है; 10 से भाग देने पर यह
$100\,000\,000$ है: दस करोड़ फेरे।
\end{solution}

\begin{solution}{exo:g5:water-shapes-the-land:11}
एक संभव कहानी। कण किसी \omterm{def:g1:landscapes-around-us:landform}{पर्वत} पर ऊँचे पर एक \omterm{def:g3:stones-pebbles-sand:rock}{चट्टान} का हिस्सा है। पाला रात-दर-रात
उसके चारों ओर एक दरार खोलता है, जब तक \omterm{def:g3:stones-pebbles-sand:rock}{चट्टान} का एक टुकड़ा टूटकर अलग नहीं हो
जाता: यह \omterm{def:g5:water-shapes-the-land:weathering}{अपक्षय} है। टुकड़ा नीचे की धारा में गिरता है। ऊपरी मार्ग में पानी तेज़ है
और ढलान खड़ी: वह टुकड़े को लुढ़काता है (\omterm{def:g5:water-shapes-the-land:erosion}{अपरदन}), और दूसरे पत्थरों से टकराते-टकराते
वह \omterm{def:g3:stones-pebbles-sand:pebble}{कंकड़ों} में और फिर \omterm{def:g3:stones-pebbles-sand:sand}{रेत} के कणों में टूट जाता है। मध्य मार्ग में \omterm{def:g1:landscapes-around-us:water}{नदी} धीमी होती
है; कण कुछ समय के लिए एक \omterm{def:g5:water-shapes-the-land:meander}{विसर्प} के मोड़ के भीतर के रेतीले \omterm{def:g1:landscapes-around-us:water}{तट} पर आराम करता है,
जब तक कोई \omterm{def:g5:water-shapes-the-land:flood}{बाढ़} उसे फिर नहीं उठा लेती। निचले मार्ग में \omterm{def:g1:landscapes-around-us:water}{नदी} धीमी और चौड़ी है; कण
पानी के साथ \omterm{def:g5:water-shapes-the-land:source-mouth}{मुहाने} तक जाता है। वहाँ पानी लगभग रुक जाता है, और कण \omterm{def:g5:water-shapes-the-land:delta}{डेल्टा} के
\omterm{def:g3:stones-pebbles-sand:sand}{कीचड़} और \omterm{def:g3:stones-pebbles-sand:sand}{रेत} के साथ बैठ जाता है: यह \omterm{def:g5:water-shapes-the-land:deposition}{निक्षेपण} है।
\end{solution}
